% GENERATED FILE. Edit canonical lessons, then run npm run generate:books.
% canonical-source-hash: fnv1a64:7a2f10c7b6cf68f5
% canonical-lessons: UR-C359-bavujud, UR-C359-alava, UR-C359-agarche, UR-C359-albatta, UR-C359-jhukna

\chapter{Even So}
\label{ch:ur-even-so}
\hlchaptermodality{359}

\begin{chapteropening}
\textbf{By the end of this chapter:} \emph{I can say what happened in spite of something.}

Complete the last lesson of chapter 359: I can say what happened in spite of something.
\end{chapteropening}

\section[\texorpdfstring{\ur{باوجود} (bāvujūd)}{باوجود (bāvujūd)}]{\ur{باوجود} (bāvujūd) — despite (\ur{کے} \ur{باوجود})}

\label{lesson:UR-C359-bavujud}



\begin{quote}
Before the new one: say the Urdu for a majority, then the Urdu for to spurn.
\end{quote}

\subsection*{You'll want to know: \ur{باوجود}}
\textbf{\ur{باوجود}} — \emph{bāvujūd} — \textquotedblleft{}despite (\ur{کے} \ur{باوجود})\textquotedblright{}.

\begin{quote}
\ur{بارش} \ur{کے} \ur{باوجود} \ur{وہ} \ur{دفتر} \ur{گیا۔}
\end{quote}

\begin{quote}
\emph{bārish ke bāvujūd voh daftar gayā.}
\end{quote}

\begin{quote}
He went to the office despite the rain.
\end{quote}

\begin{cousinweb}
\textbf{\ur{باوجود}} (\emph{bāvujūd}) is usually traced to Persian \emph{bāvujūd}, \textquotedblleft{}with the existence of\textquotedblright{}: \emph{bā}, \textquotedblleft{}with\textquotedblright{}, and \emph{vujūd}, \textquotedblleft{}existence\textquotedblright{}, a word Persian took from Arabic.
\end{cousinweb}

\begin{grammarlens}[title={What you've built}]
One more word for giving reasons and opinions.
\end{grammarlens}

\subsection*{Your turn}
\begin{itemize}
\raggedright
  \item \emph{Say it:} \emph{bāvujūd}
  \item \emph{Say it:} \emph{bāvujūd}, once more
  \item \emph{Say it:} \emph{bārish ke bāvujūd voh daftar gayā.} — He went to the office despite the rain.
  \item \emph{Recall:} say the Urdu for a majority, then the Urdu for to spurn, then say \emph{bāvujūd} again
\end{itemize}

\begin{tcolorbox}[breakable,colback=teal!4,colframe=teal!35!black,title={Before you move on}]
What does \ur{باوجود} mean? (\textquotedblleft{}Despite (\ur{کے} \ur{باوجود})\textquotedblright{}.) Say it once more.
\end{tcolorbox}
\section[\texorpdfstring{\ur{علاوہ} (alāvā)}{علاوہ (alāvā)}]{\ur{علاوہ} (alāvā) — besides; apart from (\ur{کے} \ur{علاوہ})}

\label{lesson:UR-C359-alava}



\begin{quote}
Before the new one: say the Urdu for to spurn, then the Urdu for despite.
\end{quote}

\subsection*{You'll want to know: \ur{علاوہ}}
\textbf{\ur{علاوہ}} — \emph{alāvā} — \textquotedblleft{}besides; apart from (\ur{کے} \ur{علاوہ})\textquotedblright{}.

\begin{quote}
\ur{اس} \ur{کے} \ur{علاوہ} \ur{کوئی} \ur{راستہ} \ur{نہیں} \ur{ہے۔}
\end{quote}

\begin{quote}
\emph{is ke alāvā koī rāstā nahīṅ hai.}
\end{quote}

\begin{quote}
Apart from this there is no way.
\end{quote}

\begin{cousinweb}
\textbf{\ur{علاوہ}} (\emph{alāvā}) is usually traced, through Persian, to Arabic \emph{ʿalāwa}, \textquotedblleft{}an addition\textquotedblright{}.
\end{cousinweb}

\begin{grammarlens}[title={What you've built}]
One more word for giving reasons and opinions.
\end{grammarlens}

\subsection*{Your turn}
\begin{itemize}
\raggedright
  \item \emph{Say it:} \emph{alāvā}
  \item \emph{Say it:} \emph{alāvā}, once more
  \item \emph{Say it:} \emph{is ke alāvā koī rāstā nahīṅ hai.} — Apart from this there is no way.
  \item \emph{Recall:} say the Urdu for to spurn, then the Urdu for despite, then say \emph{alāvā} again
\end{itemize}

\begin{tcolorbox}[breakable,colback=teal!4,colframe=teal!35!black,title={Before you move on}]
What does \ur{علاوہ} mean? (\textquotedblleft{}Besides; apart from (\ur{کے} \ur{علاوہ})\textquotedblright{}.) Say it once more.
\end{tcolorbox}
\section[\texorpdfstring{\ur{اگرچہ} (agarche)}{اگرچہ (agarche)}]{\ur{اگرچہ} (agarche) — even though}

\label{lesson:UR-C359-agarche}



\begin{quote}
Before the new one: say the Urdu for despite, then the Urdu for besides.
\end{quote}

\subsection*{You'll want to know: \ur{اگرچہ}}
\textbf{\ur{اگرچہ}} — \emph{agarche} — \textquotedblleft{}even though\textquotedblright{}.

\begin{quote}
\ur{اگرچہ} \ur{وہ} \ur{بیمار} \ur{تھا،} \ur{پھر} \ur{بھی} \ur{کام} \ur{پر} \ur{آیا۔}
\end{quote}

\begin{quote}
\emph{agarche voh bīmār thā, phir bhī kām par āyā.}
\end{quote}

\begin{quote}
Even though he was ill, he still came to work.
\end{quote}

\begin{cousinweb}
\textbf{\ur{اگرچہ}} (\emph{agarche}) is usually traced to Persian \emph{agarchi}, \textquotedblleft{}even if\textquotedblright{}, built on \textbf{\ur{اگر}}, \textquotedblleft{}if\textquotedblright{}.
\end{cousinweb}

\begin{grammarlens}[title={What you've built}]
One more word for giving reasons and opinions.
\end{grammarlens}

\subsection*{Your turn}
\begin{itemize}
\raggedright
  \item \emph{Say it:} \emph{agarche}
  \item \emph{Say it:} \emph{agarche}, once more
  \item \emph{Say it:} \emph{agarche voh bīmār thā, phir bhī kām par āyā.} — Even though he was ill, he still came to work.
  \item \emph{Recall:} say the Urdu for despite, then the Urdu for besides, then say \emph{agarche} again
\end{itemize}

\begin{tcolorbox}[breakable,colback=teal!4,colframe=teal!35!black,title={Before you move on}]
What does \ur{اگرچہ} mean? (\textquotedblleft{}Even though\textquotedblright{}.) Say it once more.
\end{tcolorbox}
\section[\texorpdfstring{\ur{البتہ} (albattā)}{البتہ (albattā)}]{\ur{البتہ} (albattā) — however; mind you}

\label{lesson:UR-C359-albatta}



\begin{quote}
Before the new one: say the Urdu for besides, then the Urdu for even though.
\end{quote}

\subsection*{You'll want to know: \ur{البتہ}}
\textbf{\ur{البتہ}} — \emph{albattā} — \textquotedblleft{}however; mind you\textquotedblright{}.

\begin{quote}
\ur{کام} \ur{مشکل} \ur{ہے،} \ur{البتہ} \ur{ناممکن} \ur{نہیں۔}
\end{quote}

\begin{quote}
\emph{kām mushkil hai, albattā nāmumkin nahīṅ.}
\end{quote}

\begin{quote}
The work is hard; it is not impossible, though.
\end{quote}

\begin{cousinweb}
\textbf{\ur{البتہ}} (\emph{albattā}) is usually traced, through Persian, to Arabic \emph{al-battata}, \textquotedblleft{}decidedly, by all means\textquotedblright{}; in Urdu it brings in a reservation.
\end{cousinweb}

\begin{grammarlens}[title={What you've built}]
One more word for giving reasons and opinions.
\end{grammarlens}

\subsection*{Your turn}
\begin{itemize}
\raggedright
  \item \emph{Say it:} \emph{albattā}
  \item \emph{Say it:} \emph{albattā}, once more
  \item \emph{Say it:} \emph{kām mushkil hai, albattā nāmumkin nahīṅ.} — The work is hard; it is not impossible, though.
  \item \emph{Recall:} say the Urdu for besides, then the Urdu for even though, then say \emph{albattā} again
\end{itemize}

\begin{tcolorbox}[breakable,colback=teal!4,colframe=teal!35!black,title={Before you move on}]
What does \ur{البتہ} mean? (\textquotedblleft{}However; mind you\textquotedblright{}.) Say it once more.
\end{tcolorbox}
\section[\texorpdfstring{\ur{جھکنا} (jhuknā)}{جھکنا (jhuknā)}]{\ur{جھکنا} (jhuknā) — to give in; to bow}

\label{lesson:UR-C359-jhukna}



\begin{quote}
Before the new one: say the Urdu for even though, then the Urdu for however.
\end{quote}

\subsection*{You'll want to know: \ur{جھکنا}}
\textbf{\ur{جھکنا}} — \emph{jhuknā} — \textquotedblleft{}to give in; to bow\textquotedblright{}.

\begin{quote}
\ur{مزدور} \ur{حکومت} \ur{کے} \ur{سامنے} \ur{نہیں} \ur{جھکے۔}
\end{quote}

\begin{quote}
\emph{mazdūr hukūmat ke sāmne nahīṅ jhuke.}
\end{quote}

\begin{quote}
The workers did not give in to the government.
\end{quote}

Past (he …): \ur{وہ} \ur{آخر} \ur{میں} \ur{جھک} \ur{گیا}, \emph{voh ākhir meṅ jhuk gayā}.

\begin{cousinweb}
\textbf{\ur{جھکنا}} (\emph{jhuknā}) is inherited: CDIAL lists the Hindi cousin \emph{jhuknā} under \emph{*jhukkati}, \textquotedblleft{}stoops\textquotedblright{}, a form it reconstructs (CDIAL 5399).
\end{cousinweb}

\begin{grammarlens}[title={What you've built}]
One more word for giving reasons and opinions.
\end{grammarlens}

\subsection*{Your turn}
\begin{itemize}
\raggedright
  \item \emph{Say it:} \emph{jhuknā}
  \item \emph{Say it:} \emph{jhuknā}, once more
  \item \emph{Say it:} \emph{mazdūr hukūmat ke sāmne nahīṅ jhuke.} — The workers did not give in to the government.
  \item \emph{Recall:} say the Urdu for even though, then the Urdu for however, then say \emph{jhuknā} again
\end{itemize}

\begin{tcolorbox}[breakable,colback=teal!4,colframe=teal!35!black,title={Before you move on}]
What does \ur{جھکنا} mean? (\textquotedblleft{}To give in; to bow\textquotedblright{}.) Say it once more.
\end{tcolorbox}
